\scanpage{28}
SYSTÈMES de RACINES (5) \hfill Tribu n° 59

racine simple $v_i$ envoie $v_i$ en $-v_i$ et permute les autres racines positives. Donner un lemme sur les applic. $R\longrightarrow\text{groupes}$.

Les n° 6 et 7 sont permutés.

\uline{n° 6} - (Ensembles clos de racines ; ancien n° 7).

Ajouter la prop. 4 de Serre, (n° 400, Annexe II), et sa prop. 5 ; par contre, son cor. et le lemme 2, (p. A.11), passent en exer.

Pour la prop. 15, on prend la démonstration suivante, (fournie par Borel) :

\uline{p. 110}. \uline{Prop. 15}. Démonstration :

1) Remarquons d'abord que si $x=a.y$, $(x,y\in P)$, alors $a>0$. Sinon, en effet comme $a\ne-1$ vu l'hypothèse $P\cap-P=\varnothing$ on peut, en changeant éventuellement les rôles de $x$ et $y$, supposer $a=-1/2$, d'où $x+y=y/2\in P$ et $y/2$, $-y/2\in P$, ce qui est absurde. Il s'ensuit, compte tenu de la Prop. 4 du N° 3 que si $x$, $y\in P$ et $(x,y)<0$, alors $x+y\in P$.

2) Montrons maintenant qu'aucune somme $y_1+\ldots+y_q$, $(q\geq1)$, d'éléments de $P$ n'est nulle. On procède par récurrence sur $q$ ; le cas $q=2$ résultant de $P\cap(-P)=\varnothing$, on peut supposer $q\geq3$. Si $y_1+y_q=0$, alors $-y_1=y_2+\ldots+y_q$ donc $(-y_1,y_2+\ldots+y_q)>0$, d'où l'existence d'un indice $j$ $(1\leq j\leq q)$ tel que $(y_1,y_j)<0$. On a alors $y_1+y_j\in P$ et :
\[
 (y_1+y_j)+\sum_{1\leq i\leq q,\,i\ne j}y_{\unsure{1|i}}=0,
\]
contrairement à l'hypothèse de récurrence.

3) Nous voulons montrer maintenant l'existence d'un élément $z\ne0$ de $V$ tel que $(z,x)\geq0$ quel que soit $x\in P$. Si ce n'est pas le cas, alors, vu (1), on

\mspage{27}
